% GENERATED FILE. Edit canonical lessons, then run npm run generate:books.
% canonical-source-hash: fnv1a64:2bc849c736bb7a65
% canonical-lessons: ES-C685-paisaje, ES-C685-naturaleza, ES-C685-sierra, ES-C685-monte, ES-C685-desierto

\chapter{El paisaje}
\label{ch:es-el-paisaje}
\hlchaptermodality{685}

\begin{chapteropening}
\textbf{By the end of this chapter:} \emph{I can describe a landscape: nature, mountain ranges, hills and deserts.}

Complete the last lesson of chapter 685: I can describe a landscape: nature, mountain ranges, hills and deserts.
\end{chapteropening}

\section[el paisaje]{el paisaje — the landscape, the scenery}

\label{lesson:ES-C685-paisaje}



\begin{quote}
Before the new one: say the Spanish for to reside, then the Spanish for the neighbourhood.
\end{quote}

\subsection*{You'll want to know: el paisaje}
\textbf{el paisaje} — \textquotedblleft{}the landscape, the scenery\textquotedblright{}.

\begin{quote}
El paisaje es precioso.
\end{quote}

\begin{quote}
\emph{The scenery is beautiful.}
\end{quote}

\emph{Un paisaje de montaña} — a mountain landscape.

\begin{cousinweb}
\emph{paisaje} was borrowed from French \emph{paysage}, built on \emph{pays} \textquotedblleft{}country\textquotedblright{}, which goes back to Latin \emph{pāgus} \textquotedblleft{}a district, the countryside\textquotedblright{}.
\end{cousinweb}

\begin{grammarlens}[title={What you've built}]
One more word for telling, explaining and describing at B1.
\end{grammarlens}

\subsection*{Your turn}
\begin{itemize}
\raggedright
  \item \emph{Say it:} \emph{el paisaje}
  \item \emph{Say it:} \emph{el paisaje}, once more
  \item \emph{Say it:} a sentence of your own with \emph{el paisaje}
  \item \emph{Recall:} say the Spanish for to reside, then the Spanish for the neighbourhood, then say \emph{el paisaje} again
\end{itemize}

\begin{tcolorbox}[breakable,colback=teal!4,colframe=teal!35!black,title={Before you move on}]
What does el paisaje mean? (\textquotedblleft{}The landscape, the scenery\textquotedblright{}.) Say it once more.
\end{tcolorbox}
\section[la naturaleza]{la naturaleza — nature}

\label{lesson:ES-C685-naturaleza}



\begin{quote}
Before the new one: say the Spanish for the neighbourhood, then the Spanish for the landscape.
\end{quote}

\subsection*{You'll want to know: la naturaleza}
\textbf{la naturaleza} — \textquotedblleft{}nature\textquotedblright{}.

\begin{quote}
Me gusta pasear por la naturaleza.
\end{quote}

\begin{quote}
\emph{I like walking in nature.}
\end{quote}

\emph{Por naturaleza} — by nature.

\begin{cousinweb}
\emph{naturaleza} is built on \emph{natural} with \emph{-eza} (EtymDB); \emph{natural} comes from Latin \emph{nātūrālis}, built on \emph{nātus} \textquotedblleft{}born\textquotedblright{}.
\end{cousinweb}

\begin{grammarlens}[title={What you've built}]
One more word for telling, explaining and describing at B1.
\end{grammarlens}

\subsection*{Your turn}
\begin{itemize}
\raggedright
  \item \emph{Say it:} \emph{la naturaleza}
  \item \emph{Say it:} \emph{la naturaleza}, once more
  \item \emph{Say it:} a sentence of your own with \emph{la naturaleza}
  \item \emph{Recall:} say the Spanish for the neighbourhood, then the Spanish for the landscape, then say \emph{la naturaleza} again
\end{itemize}

\begin{tcolorbox}[breakable,colback=teal!4,colframe=teal!35!black,title={Before you move on}]
What does la naturaleza mean? (\textquotedblleft{}Nature\textquotedblright{}.) Say it once more.
\end{tcolorbox}
\section[la sierra]{la sierra — the mountain range; the saw}

\label{lesson:ES-C685-sierra}



\begin{quote}
Before the new one: say the Spanish for the landscape, then the Spanish for nature.
\end{quote}

\subsection*{You'll want to know: la sierra}
\textbf{la sierra} — \textquotedblleft{}the mountain range; the saw\textquotedblright{}.

\begin{quote}
Pasamos el fin de semana en la sierra.
\end{quote}

\begin{quote}
\emph{We spent the weekend in the mountains.}
\end{quote}

A row of peaks looks like the teeth of a saw.

\begin{cousinweb}
\emph{sierra} is inherited from Latin \emph{serra} \textquotedblleft{}a saw\textquotedblright{} (EtymDB); a jagged line of peaks looks like its teeth.
\end{cousinweb}

\begin{grammarlens}[title={What you've built}]
One more word for telling, explaining and describing at B1.
\end{grammarlens}

\subsection*{Your turn}
\begin{itemize}
\raggedright
  \item \emph{Say it:} \emph{la sierra}
  \item \emph{Say it:} \emph{la sierra}, once more
  \item \emph{Say it:} a sentence of your own with \emph{la sierra}
  \item \emph{Recall:} say the Spanish for the landscape, then the Spanish for nature, then say \emph{la sierra} again
\end{itemize}

\begin{tcolorbox}[breakable,colback=teal!4,colframe=teal!35!black,title={Before you move on}]
What does la sierra mean? (\textquotedblleft{}The mountain range; the saw\textquotedblright{}.) Say it once more.
\end{tcolorbox}
\section[el monte]{el monte — the mountain, the hill; the woodland}

\label{lesson:ES-C685-monte}



\begin{quote}
Before the new one: say the Spanish for nature, then the Spanish for the mountain range.
\end{quote}

\subsection*{You'll want to know: el monte}
\textbf{el monte} — \textquotedblleft{}the mountain, the hill; the woodland\textquotedblright{}.

\begin{quote}
Subimos al monte por la mañana.
\end{quote}

\begin{quote}
\emph{We went up the mountain in the morning.}
\end{quote}

\emph{El monte} is also wild land with trees and bushes.

\begin{cousinweb}
\emph{monte} is inherited from Latin \emph{montem}, a form of \emph{mōns} \textquotedblleft{}a mountain\textquotedblright{} (EtymDB).
\end{cousinweb}

\begin{grammarlens}[title={What you've built}]
One more word for telling, explaining and describing at B1.
\end{grammarlens}

\subsection*{Your turn}
\begin{itemize}
\raggedright
  \item \emph{Say it:} \emph{el monte}
  \item \emph{Say it:} \emph{el monte}, once more
  \item \emph{Say it:} a sentence of your own with \emph{el monte}
  \item \emph{Recall:} say the Spanish for nature, then the Spanish for the mountain range, then say \emph{el monte} again
\end{itemize}

\begin{tcolorbox}[breakable,colback=teal!4,colframe=teal!35!black,title={Before you move on}]
What does el monte mean? (\textquotedblleft{}The mountain, the hill; the woodland\textquotedblright{}.) Say it once more.
\end{tcolorbox}
\section[el desierto]{el desierto — the desert; deserted}

\label{lesson:ES-C685-desierto}



\begin{quote}
Before the new one: say the Spanish for the mountain range, then the Spanish for the mountain.
\end{quote}

\subsection*{You'll want to know: el desierto}
\textbf{el desierto} — \textquotedblleft{}the desert; deserted\textquotedblright{}.

\begin{quote}
Cruzamos el desierto en coche.
\end{quote}

\begin{quote}
\emph{We drove across the desert.}
\end{quote}

As an adjective it means empty: \emph{una calle desierta}.

\begin{cousinweb}
\emph{desierto} is inherited from Latin \emph{dēsertus} \textquotedblleft{}abandoned, forsaken\textquotedblright{} (EtymDB): a desert is land that people have left.
\end{cousinweb}

\begin{grammarlens}[title={What you've built}]
One more word for telling, explaining and describing at B1.
\end{grammarlens}

\subsection*{Your turn}
\begin{itemize}
\raggedright
  \item \emph{Say it:} \emph{el desierto}
  \item \emph{Say it:} \emph{el desierto}, once more
  \item \emph{Say it:} a sentence of your own with \emph{el desierto}
  \item \emph{Recall:} say the Spanish for the mountain range, then the Spanish for the mountain, then say \emph{el desierto} again
\end{itemize}

\begin{tcolorbox}[breakable,colback=teal!4,colframe=teal!35!black,title={Before you move on}]
What does el desierto mean? (\textquotedblleft{}The desert; deserted\textquotedblright{}.) Say it once more.
\end{tcolorbox}
